\documentclass[11pt]{article}
\usepackage[margin=1in]{geometry}
\usepackage{amsmath, amssymb}

\setlength{\parindent}{0pt}
\setlength{\parskip}{0.6em}

\begin{document}

%------------------------------------------------------------
\section*{1.)}

If $x^2 + 5x$ is odd then $x$ is odd.

Contrapositive: if $x$ is even then $x^2 + 5x$ is even.

A number $n$ is even iff $2k = n$.

Since we're assuming $x$ is even, we can write it as $2k$.

Likewise with $x^2 + 5x \Rightarrow (2k)^2 + 5(2k) \to 4k^2 + 10k$

$2\underbrace{(2k^2 + 5k)}_{\text{int.}} \leftarrow$ Definitionally even.

We assumed $x$ is even $(2k)$, and $x^2 + 5x$ also even $(2k)^2 + 5(2k)$ which proves true.

\hfill $\square$ \quad Since the contrapositive is true, the original claim is true.

\bigskip\hrule\bigskip

%------------------------------------------------------------
\section*{2.)}

Given $x_1 < x_2$, if $\underbrace{f(x_1) > f(x_2)}_{p}$ then $\underbrace{f(x) \text{ is not increasing}}_{q}$.

$\sim q \Rightarrow \sim p =$ Given $x_1 < x_2$, if $f(x)$ is increasing then $f(x_1) < f(x_2)$.

$\bigstar$ $f(x)$ is increasing on domain $\mathbb{R}$ iff $\forall x \in \mathbb{R},\ \forall k > 0,\ f(x) < f(x + k)$.

Let $x_2 - x_1 = k$. $k$ here is $> 0$ since $x_1 < x_2$.

Since $k = x_2 - x_1$, we can rewrite it as

$x_2 = x_1 + k$

Using the definition, substituting $x = x_1$ and $k = x_2 - x_1$, so $x + k = x_2$

$f(x_1) < f(x_2)$

\hfill $\square$ \quad Therefore, given $x_1 < x_2$, when $f(x)$ is increasing then $f(x_1) < f(x_2)$.

\bigskip\hrule\bigskip

%------------------------------------------------------------
\section*{3.)}

If $\underbrace{f(x) = x^2}_{p}$ then $\underbrace{f(x) \text{ is not bounded above}}_{q}$.

$p \to {\sim}q =$ if $f(x) = x^2$ then $f(x)$ is bounded above.

$\bigstar$ $f(x)$ is bounded above iff $\exists K \in \mathbb{R},\ \forall x \in \mathbb{R},\ f(x) \leq K$


Since $f(x)$ is $x^2$ in this case, we can sub $x^2$ for $f(x)$ in our definition:

$\exists K \in \mathbb{R},\ \forall x \in \mathbb{R},\ x^2 \leq K$.

Let $x = |K| + 1$ Since that guarantees $x > K$ always.

Sub $|K| + 1$ for $x$ in $x^2$: $(|K| + 1)^2 \to K^2 + 2|K| + 1$

  \[
  \left.
  \begin{aligned}
  &K^2 \text{ is always } 0 \text{ or } > 0\\
  &2|K| \text{ is always } 0 \text{ or } > 0\\
  &1 \text{ is always } > 0
  \end{aligned}
  \right\} \text{all always greater than } K.
  \]


\hfill $\square$ \quad We have found an $x$ $(|K| + 1)$ where $x^2 > K$, which contradicts our claim. Therefore, the original claim must be true.

\bigskip\hrule\bigskip

%------------------------------------------------------------
\section*{4.)}

If $\underbrace{3 \mid (n^2 + 2)}_{p}$ $\to$ $\underbrace{3 \nmid n}_{q}$

$p \to {\sim}q =$ If $3 \mid (n^2 + 2)$ then $3 \mid n$

$\bigstar$ $a \mid b$ iff $\exists c \in \mathbb{Z},\ ac = b$.

Using our definition, $3 \mid n = n = 3c$

Similarly, $3 \mid (n^2 + 2) = n^2 + 2 = 3d$

Since $n = 3c$, we can substitute $3c$ for $n$ in $n^2 + 2$

$(3c)^2 + 2 = 3d \to 9c^2 + 2 = 3d \to 2 = 3d - 9c^2$

$2 = 3\underbrace{(d - 3c^2)}_{m} = \underline{\underline{2 = 3m}} =$ by definition is $3 \mid 2$

Which is false because $\frac{2}{3} \notin \mathbb{Z}$

\hfill $\square$ \quad Therefore the original claim must be true.

\end{document}